% In this file you should put the actual content of the blueprint.
% It will be used both by the web and the print version.
% It should *not* include the \begin{document}
%
% If you want to split the blueprint content into several files then
% the current file can be a simple sequence of \input. Otherwise It
% can start with a \section or \chapter for instance.

\section*{Introduction}

\section{The unique choice axiom}

\section{Basic Definitions}
Let $R$ be a commutative ring.

\begin{definition}
  \label{SDG.D}
  \leanok
  \lean{SDG.D}
We denote by $D$ the subsemigroup of $R$
\[
\{r \in R \mbox{ such that } r ^ 2 = 0 \}
\]
\end{definition}

\begin{definition}
  \label{SDG.𝔻}
  \leanok
  \lean{SDG.𝔻}
For any $k\in \mathbb{N},$ we denote by $D_k$ the subsemigroup of $R$
\[
\{r \in R \mbox{ such that } r ^ {k+1} = 0 \}
\]
\end{definition}

\begin{lemma}
  \label{SDG.zero_mem_𝔻}
  \leanok
  \lean{SDG.zero_mem_D}
  \uses{SDG.D}
  We have that $0 \in D.$
\end{lemma}
\begin{proof}
  \leanok
  Obvious.
\end{proof}

\begin{definition}
  \label{SDG.IsKockLawvereone}
  \leanok
  \lean{SDG.IsKockLawvereone}
  \uses{SDG.D}
We say that $R$ is \emph{Kock--Lawvere} if it is nontrivial and forall $g : D → R$ there is a unique
$b \in R$ such that forall $d \in D$ we have $g(d) = g(0) + d \cdot b$.
\end{definition}

\begin{definition}
  \label{SDG.deriv}
  \leanok
  \lean{SDG.deriv}
  \uses{SDG.IsKockLawvereone}
  Let $f \colon R \to R$ be any function. We define $\partial f \colon R \to R$, the \emph{derivative} of $f$, as the unique function such that for all $x \in R$ and $d \in D$ we have
\[
f(x + d)=f(0)+\partial f(x) \cdot d
\]
We have that $\partial$ is a $R$-linear operator and moreover it satisfies the \emph{Leibnitz rule}: for all function $f$ and $g$ we have
\[
\partial(f \cdot g) = \partial f \cdot g + f \cdot \partial g
\]
\end{definition}

\section{The synthetic theory}
